\section{结论}

本文针对AUV海流估计中水动力模型失配的鲁棒性问题，提出了一种激励门控CFD先验海流观测器(EG-UCCO)。主要贡献包括：

\begin{enumerate}
    \item \textbf{方法论创新}：利用CFD预标定阻力模型作为确定性先验，设计速度层面的预测-校正观测器架构，避免了加速度差分导致的噪声放大问题。灵敏度Gramian激励门控基于Fisher信息理论，比经验门控具有更严格的数学基础。

    \item \textbf{系统性验证}：在3种海流场景、4种模型失配水平下进行了全面的对比实验。EG-UCCO在所有12个配置中均取得最优海流估计精度，在30\% CFD阻力系数扰动下估计精度仅退化4\%，优于调优EKF的16\%退化。

    \item \textbf{机制证据}：门控消融实验证实Gramian门控在传感器噪声下可提升23\%的估计精度，且门控收益随噪声增大而增加，验证了激励门控作为鲁棒性机制的有效性。

    \item \textbf{工程洞察}：补偿消融实验揭示了XHY AUV的Surge-Yaw通道耦合现象，为前馈补偿的设计选择（仅Yaw补偿）提供了实验依据。
\end{enumerate}

未来工作将沿以下方向展开：(1) 在实际海洋环境中进行海试验证，以独立ADCP测量作为海流真值参考；(2) 将EG-UCCO扩展到多AUV协同海流估计，利用编队信息提升空间流场的估计精度；(3) 研究CFD模型在线更新机制，应对生物附着等因素导致的长期模型漂移。
